\documentclass[11pt]{article}

\usepackage{amsmath,amssymb}
\usepackage{xurl}
\usepackage[colorlinks=true,linkcolor=blue,citecolor=blue,urlcolor=blue]{hyperref}
\setlength{\emergencystretch}{2em}

\input{./command.tex}

\title{Yamagami's Simultaneous-Singularity Boundary Count}
\author{}
\date{2026-08-24}

\begin{document}
\maketitle
\gapnote{clarification}{resolved}

\section{The cyclic function and its boundary}

In the case denoted by $l=1$, Yamagami studies a cyclic function on
$\mathbb Z/N\mathbb Z$ \cite[pp.~524--525]{Yamagami1993Cyclic}. For a vector
$x=(x_q)_{q\in\mathbb Z/N\mathbb Z}$, define
\begin{align}
    f_{N,m,s}(x)
        &=
        \sum_{q\in\mathbb Z/N\mathbb Z}
        \frac{x_q}
        {(s-m)x_q+x_{q+1}+\cdots+x_{q+m}}.
    \label{eq:yamagami93_simultaneous_singularity_boundary_cyclic_function}
\end{align}
All indices in
(\ref{eq:yamagami93_simultaneous_singularity_boundary_cyclic_function}) are
cyclic. Throughout the boundary count, assume
\begin{align}
    N
        &\geq 1,
    \label{eq:yamagami93_simultaneous_singularity_boundary_dimension}\\
    m
        &< N,
    \label{eq:yamagami93_simultaneous_singularity_boundary_window_length}\\
    N
        &\leq s.
    \label{eq:yamagami93_simultaneous_singularity_boundary_parameter_bound}
\end{align}
These assumptions imply
\begin{align}
    m
        &< s,
    \label{eq:yamagami93_simultaneous_singularity_boundary_strict_parameter}\\
    s-m
        &> 0.
    \label{eq:yamagami93_simultaneous_singularity_boundary_positive_diagonal}
\end{align}

Let $\{x^{(r)}\}_{r\geq 1}$ be a sequence of strictly positive vectors that
converges coordinatewise to a nonzero boundary vector
$a=(a_q)_{q\in\mathbb Z/N\mathbb Z}$ with $a_q\geq 0$. If the limiting
denominator at an index $i$ vanishes, then
\begin{align}
    (s-m)a_i+a_{i+1}+\cdots+a_{i+m}
        &= 0.
    \label{eq:yamagami93_simultaneous_singularity_boundary_zero_denominator}
\end{align}
Every term on the left-hand side of
(\ref{eq:yamagami93_simultaneous_singularity_boundary_zero_denominator}) is
nonnegative, and
(\ref{eq:yamagami93_simultaneous_singularity_boundary_positive_diagonal})
makes the coefficient of $a_i$ positive. Hence
\begin{align}
    a_i
        &= a_{i+1}
         = \cdots
         = a_{i+m}
         = 0.
    \label{eq:yamagami93_simultaneous_singularity_boundary_zero_window}
\end{align}
Because $a\neq 0$, cyclic traversal after this forward zero window eventually
reaches a positive coordinate.

\section{The calculation printed in the source}

Yamagami first treats one indeterminate summand
\cite[p.~525]{Yamagami1993Cyclic}. After a cyclic relabeling, the hypotheses
are
\begin{align}
    a_N
        &= a_1
         = \cdots
         = a_m
         = 0,
    \label{eq:yamagami93_simultaneous_singularity_boundary_single_zero_run}\\
    a_{m+1}
        &> 0,
    \label{eq:yamagami93_simultaneous_singularity_boundary_single_right_positive}\\
    a_{N-1}
        &> 0.
    \label{eq:yamagami93_simultaneous_singularity_boundary_single_left_positive}
\end{align}
The $m$ summands with numerators indexed by $1,\ldots,m$ tend to zero. Every
other positive summand is at most $1/(s-m)$. The printed estimates therefore
give
\begin{align}
    \limsup_{r\to\infty} f_{N,m,s}(x^{(r)})
        &\leq \frac{N-m}{s-m},
    \label{eq:yamagami93_simultaneous_singularity_boundary_single_limsup}\\
    \frac{N-m}{s-m}
        &\leq \frac{N}{s}.
    \label{eq:yamagami93_simultaneous_singularity_boundary_ratio_comparison}
\end{align}

For several indeterminate summands, Ref.~\cite{Yamagami1993Cyclic} records a
cyclic zero-string condition. It selects an index $j$ such that
\begin{align}
    a_j
        &> 0,
    \label{eq:yamagami93_simultaneous_singularity_boundary_run_endpoint}\\
    a_{j-m-1}
        &= a_{j-m}
         = \cdots
         = a_{j-1}
         = 0.
    \label{eq:yamagami93_simultaneous_singularity_boundary_singular_run}
\end{align}
The paper says that the remaining argument is similar, but it does not print
the finite first-positive construction or the simultaneous limsup
calculation. Thus
(\ref{eq:yamagami93_simultaneous_singularity_boundary_single_zero_run})--%
(\ref{eq:yamagami93_simultaneous_singularity_boundary_ratio_comparison})
constitute the printed single-singularity calculation, while the reduction of
an arbitrary family of singular denominators to one global count is omitted.

\section{The completed simultaneous count}

Begin with the zero window in
(\ref{eq:yamagami93_simultaneous_singularity_boundary_zero_window}). Traverse
the cyclic coordinates after $i+m$, and choose the first positive coordinate.
Finiteness of $\mathbb Z/N\mathbb Z$ and the assumption $a\neq 0$ ensure that
such a coordinate $a_j$ exists. Minimality implies that every coordinate from
the end of the zero window through $j-1$ vanishes. In particular, it gives
(\ref{eq:yamagami93_simultaneous_singularity_boundary_run_endpoint}) and
(\ref{eq:yamagami93_simultaneous_singularity_boundary_singular_run}).\footnote{The
finite first-positive step is formalized by
\leanid{Yamagami.exists_hasSingularRun_of_zero_window}; its application to a
zero denominator is
\leanid{Yamagami.exists_hasSingularRun_of_forwardDenominator_eq_zero} in
\doclink{QICLean/Analysis/YamagamiBoundary}.}

Define the cyclic index set
\begin{align}
    V
        &=
        \{j-1,j-2,\ldots,j-m\}.
    \label{eq:yamagami93_simultaneous_singularity_boundary_vanishing_indices}
\end{align}
The assumption
(\ref{eq:yamagami93_simultaneous_singularity_boundary_window_length}) makes
these $m$ cyclic indices distinct. Therefore,
\begin{align}
    |V|
        &= m,
    \label{eq:yamagami93_simultaneous_singularity_boundary_vanishing_cardinality}\\
    \left|(\mathbb Z/N\mathbb Z)\setminus V\right|
        &= N-m.
    \label{eq:yamagami93_simultaneous_singularity_boundary_complement_cardinality}
\end{align}
For $q=j-k\in V$, the limiting numerator $a_q$ is zero, while the forward
denominator contains the fixed positive coordinate $a_j$. The corresponding
summand therefore tends to zero. Every summand indexed outside $V$ remains
nonnegative and is bounded above by $1/(s-m)$ along the positive
approximants. Partitioning the finite sum once gives
\begin{align}
    \limsup_{r\to\infty} f_{N,m,s}(x^{(r)})
        &\leq
        \left|(\mathbb Z/N\mathbb Z)\setminus V\right|
        \frac{1}{s-m},
    \label{eq:yamagami93_simultaneous_singularity_boundary_partition_bound}\\
    \limsup_{r\to\infty} f_{N,m,s}(x^{(r)})
        &\leq
        \frac{N-m}{s-m}.
    \label{eq:yamagami93_simultaneous_singularity_boundary_global_limsup}
\end{align}
The assumption
(\ref{eq:yamagami93_simultaneous_singularity_boundary_parameter_bound}) gives
the comparison in
(\ref{eq:yamagami93_simultaneous_singularity_boundary_ratio_comparison}).\footnote{The
reusable finite limsup estimate is
\leanid{Yamagami.limsup_sum_le_card_mul_of_tendsto_zero_on}. The
boundary-condition form of
(\ref{eq:yamagami93_simultaneous_singularity_boundary_global_limsup}) and its
comparison with $N/s$ are
\leanid{Yamagami.limsup_functional_le_sub_ratio_of_singularDenominator} and
\leanid{Yamagami.limsup_functional_le_card_ratio_of_singularDenominator} in
\doclink{QICLean/Analysis/YamagamiBoundary}.}

This argument uses one global count. It neither assigns singular denominators
to maximal zero blocks nor adds estimates obtained from separate blocks.
Consequently, it applies without change when any number of limiting
denominators vanish.

\section{Scope of the resolution}

The formal development also proves a conditional passage to the direct
nonnegative value of
(\ref{eq:yamagami93_simultaneous_singularity_boundary_cyclic_function}), using
the convention $0/0=0$ for real division. It assumes the corresponding
inequality for every strictly positive vector. Given a nonnegative vector
$x$, it uses the uniform perturbation
\begin{align}
    y_q^{(r)}
        &= x_q+\frac{1}{r+1}.
    \label{eq:yamagami93_simultaneous_singularity_boundary_uniform_perturbation}
\end{align}
This conditional passage is formalized separately.\footnote{Formal
declarations:
\leanid{Yamagami.functional_le_of_strictlyPositive} and
\leanid{Yamagami.functional_le_card_ratio_of_strictlyPositive} in
\doclink{QICLean/Analysis/YamagamiBoundary}.}
It does not supply the strictly positive interior inequality.

The regular-boundary recurrence on p.~525 of
Ref.~\cite{Yamagami1993Cyclic} is formalized by
\leanid{Yamagami.functional_le_deleteLastCoordinate}. The finite-coordinate
Nowosad theorem~\cite{Nowosad1968Isoperimetric} and the uniqueness assertion
in Lemma~3 of Ref.~\cite{Yamagami1993Cyclic} are formalized by
\leanid{Nowosad.finiteCoordinate_theorem_one_eight} and
\leanid{Yamagami.lemma_three_localMax_isScalar}, respectively. The cyclic
matrix package in \doclink{QICLean/Analysis/YamagamiCyclicMatrix} verifies the
Fourier--Hessian hypotheses and identifies the generic Nowosad functional with
the concrete cyclic functional \leanid{Yamagami.functional}.

The declaration \leanid{Yamagami.functional_card_le_one} assembles these
ingredients. Its compact-superlevel argument produces a positive maximizer,
and the Fourier--Hessian and Nowosad uniqueness theorem makes the maximizer
scalar.

\section{Verdict}

The omitted simultaneous-singularity calculation is a resolved clarification.
The formal lemma proves the global counting and limsup argument suggested by
(\ref{eq:yamagami93_simultaneous_singularity_boundary_run_endpoint}) and
(\ref{eq:yamagami93_simultaneous_singularity_boundary_singular_run}) without
attributing the unprinted calculation to Ref.~\cite{Yamagami1993Cyclic}. This
resolution supplies one boundary input to the subsequently completed cyclic
inequality.

\bibliographystyle{alpha}
\bibliography{references}

\end{document}
